\chapter{Programaci\'on en Base de Datos: PL/pgSQL, Procedimientos y Triggers}
\label{cap:programacion_plpgsql_triggers}

\section{Marco Conceptual y Te\'orico (Curr\'iculo Universitario)}

La arquitectura procedimental en servidor permite ejecutar bloques de c\'odigo complejos con control transaccional, reduciendo los saltos de red (\textit{network roundtrips}) entre la aplicaci\'on y el motor de base de datos.

\subsection{Estructura de Bloques en PL/pgSQL}
Todo bloque procedimental se define mediante la sintaxis delimitada:
\begin{lstlisting}[style=sqlstyle]
[ <<label>> ]
[ DECLARE
    declarations ]
BEGIN
    statements
[ EXCEPTION
    WHEN condition THEN
        handler_statements ]
END [ label ];
\end{lstlisting}

\subsection{Disparadores (Triggers) y Eventos}
Un trigger es una funci\'on que se ejecuta autom\'aticamente cuando ocurre un evento de modificaci\'on de datos (\texttt{INSERT}, \texttt{UPDATE}, \texttt{DELETE} o \texttt{TRUNCATE}):
\begin{itemize}
    \item \textbf{Nivel de Fila (\texttt{FOR EACH ROW}):} Se ejecuta por cada registro afectado. Tiene acceso a las variables de contexto pseudo-registro \texttt{NEW} y \texttt{OLD}.
    \item \textbf{Nivel de Sentencia (\texttt{FOR EACH STATEMENT}):} Se ejecuta una sola vez por cada instrucci\'on SQL, independientemente del n\'umero de filas modificadas.
    \item \textbf{Tiempos de Disparo:} \texttt{BEFORE} (para validaciones o transformaciones previas), \texttt{AFTER} (para registros de auditor\'ia y cascadas) o \texttt{INSTEAD OF} (sobre vistas).
\end{itemize}

---

\section{Casos de Estudio y Pr\'actica Aplicada}

\begin{lstlisting}[style=sqlstyle, caption={Procedimiento almacenado con transacci\'on y validaci\'on de saldo}]
CREATE OR REPLACE PROCEDURE transfer_balance(
    p_sender_id BIGINT,
    p_receiver_id BIGINT,
    p_amount NUMERIC(12, 2)
)
LANGUAGE plpgsql
AS $$
DECLARE
    v_sender_balance NUMERIC(12, 2);
BEGIN
    -- Validar monto
    IF p_amount <= 0 THEN
        RAISE EXCEPTION 'El monto a transferir debe ser mayor a cero. Monto ingresado: %', p_amount;
    END IF;

    -- Bloqueo pesimista de cuenta emisora
    SELECT balance INTO v_sender_balance
    FROM bank_accounts
    WHERE account_id = p_sender_id
    FOR UPDATE;

    IF NOT FOUND THEN
        RAISE EXCEPTION 'Cuenta emisora ID % no encontrada', p_sender_id;
    END IF;

    IF v_sender_balance < p_amount THEN
        RAISE EXCEPTION 'Saldo insuficiente. Saldo actual: %, Requerido: %', v_sender_balance, p_amount;
    END IF;

    -- D\'ebito y Cr\'edito
    UPDATE bank_accounts SET balance = balance - p_amount WHERE account_id = p_sender_id;
    UPDATE bank_accounts SET balance = balance + p_amount WHERE account_id = p_receiver_id;

    COMMIT;
END;
$$;
\end{lstlisting}

---

\section{Taller de Consolidaci\'on del Cap\'itulo}

\begin{businessproblem}
Implementar un sistema de auditor\'ia estricto en PostgreSQL para la tabla \texttt{products} de Northwind. Cada vez que el precio unitario de un producto sea modificado, el sistema debe registrar en una tabla de auditor\'ia el ID del producto, el precio anterior, el nuevo precio, la diferencia porcentual, el usuario que ejecut\'o la acci\'on y la marca de tiempo exacta, impidiendo incrementos superiores al 50\% de un solo cambio.
\end{businessproblem}

\subsection{Soluci\'on T\'ecnica con Trigger en PostgreSQL}

\begin{lstlisting}[style=sqlstyle, caption={Trigger con validaci\'on de umbral y auditor\'ia}]
-- 1. Tabla de Registro de Auditor\'ia de Precios
CREATE TABLE product_price_audit (
    audit_id BIGSERIAL PRIMARY KEY,
    product_id SMALLINT NOT NULL,
    old_unit_price REAL NOT NULL,
    new_unit_price REAL NOT NULL,
    percentage_change NUMERIC(6, 2) NOT NULL,
    changed_by VARCHAR(80) DEFAULT CURRENT_USER,
    changed_at TIMESTAMPTZ DEFAULT CURRENT_TIMESTAMP
);

-- 2. Funci\'on de Trigger con Regla de Negocio
CREATE OR REPLACE FUNCTION trg_validate_and_audit_price()
RETURNS TRIGGER
LANGUAGE plpgsql
AS $$
DECLARE
    v_pct_change NUMERIC(6, 2);
BEGIN
    -- Solo actuar si el precio efectivamente cambi\'o
    IF NEW.unit_price <> OLD.unit_price THEN
        -- Calcular porcentaje de cambio
        v_pct_change := ((NEW.unit_price - OLD.unit_price) / OLD.unit_price) * 100.0;

        -- Regla de Negocio: Bloquear incrementos mayores al 50%
        IF v_pct_change > 50.0 THEN
            RAISE EXCEPTION 'Operaci\'on rechazada: Incremento de precio del % supera el umbral permitido del 50%% para el producto ID %',
                ROUND(v_pct_change, 2), OLD.product_id;
        END IF;

        -- Insertar en auditor\'ia
        INSERT INTO product_price_audit (
            product_id,
            old_unit_price,
            new_unit_price,
            percentage_change
        ) VALUES (
            OLD.product_id,
            OLD.unit_price,
            NEW.unit_price,
            v_pct_change
        );
    END IF;

    RETURN NEW;
END;
$$;

-- 3. Asignaci\'on del Trigger a la Tabla
CREATE TRIGGER trigger_audit_product_price
BEFORE UPDATE OF unit_price ON products
FOR EACH ROW
EXECUTE FUNCTION trg_validate_and_audit_price();
\end{lstlisting}

\begin{engtip}
\begin{itemize}
    \item \textbf{\texttt{BEFORE UPDATE OF column}:} Al especificar la columna objetivo en el evento (\texttt{OF unit\_price}), PostgreSQL optimiza el rendimiento evitando invocar la funci\'on del trigger cuando se actualizan otras columnas como el stock o la descripci\'on.
\end{itemize}
\end{engtip}
